\chapter{Conceptual Study}
\label{ch:conceptual}

\section{Introduction}
This chapter presents the analysis and design phase of the Tahkik project. We use the Unified Modeling Language (UML) to model the system's requirements, behavior, and structure. The chapter begins with requirements identification, proceeds to behavioral modeling through use case and sequence diagrams, and concludes with the static structure via class diagrams and database schema design.

\section{System Development Process}
Our development process is inspired by Agile Methodology, based on flexibility, iterative development and continuous feedback. Our workflow follows these practices:

\phaseblock{1}{Backlog Management}{We maintain a flexible list of features and tasks. Every 10 to 15 days, we select a set of prioritized tasks forming our Sprint Backlog, including implementation (mobile and web), UI/UX design, AI integration, and system conception.}

\phaseblock{2}{Sprint Cycle (1--2 Weeks)}{Each sprint focuses on producing a specific feature. The team plans upcoming tasks, implements them based on each member's role, reviews completed work, and provides feedback for improvement.}

\phaseblock{3}{Continuous Review}{We conduct continuous informal reviews where each member presents their work, and the rest of the team provides constructive feedback on quality and direction.}

\phaseblock{4}{Incremental Production}{At the end of each sprint, a part of the platform is complete and integrated, leading to a fully functional product through incremental releases.}

\section{Requirements Analysis}
\label{sec:cs_requirements}

\subsection{Functional Requirements}

\begin{table}[H]
  \centering
  \caption{Functional requirements}
  \label{tab:cs_func_req}
  \renewcommand{\arraystretch}{1.3}
  \begin{tabular}{clp{8cm}}
    \toprule
    \textbf{ID} & \textbf{Requirement} & \textbf{Description} \\
    \midrule
    FR-01 & Authentication & Users must be able to register, log in (email/password or Google Sign-In), and log out securely. \\
    FR-02 & Quran Reading & Learners can browse and read the Quran with Warsh script rendering. \\
    FR-03 & Recitation Recording & Learners can record their voice reciting selected verses. \\
    FR-04 & AI Feedback & The system processes audio using Whisper AI and returns a confidence score with error detection. \\
    FR-05 & Teacher Review & Teachers can listen to submitted recitations, view AI feedback, and provide manual corrections. \\
    FR-06 & Progress Tracking & Learners can view their recitation history, scores, and improvement trends. \\
    FR-07 & Student Management & Teachers can manage their assigned students and track individual progress. \\
    FR-08 & Subscription & The system handles different subscription tiers for premium features. \\
    FR-09 & Ijaza Certification & Teachers can issue digital Ijaza certificates to qualifying students. \\
    FR-10 & Platform Analytics & Admins can view overall platform statistics (users, recitations, AI health). \\
    \bottomrule
  \end{tabular}
\end{table}

\subsection{Non-Functional Requirements}

\begin{table}[H]
  \centering
  \caption{Non-functional requirements}
  \label{tab:cs_nonfunc_req}
  \renewcommand{\arraystretch}{1.3}
  \begin{tabular}{clp{8cm}}
    \toprule
    \textbf{ID} & \textbf{Category} & \textbf{Description} \\
    \midrule
    NFR-01 & Performance & AI feedback should return results within a few seconds of audio submission. \\
    NFR-02 & Security & User data and audio must be securely stored using JWT and RLS policies. \\
    NFR-03 & Usability & The mobile interface must be intuitive for users of all ages and technical levels. \\
    NFR-04 & Scalability & The backend must handle concurrent audio processing requests. \\
    NFR-05 & Reliability & The platform should maintain high availability with minimal downtime. \\
    \bottomrule
  \end{tabular}
\end{table}

\section{Use Case Modeling}
\label{sec:cs_usecases}

\begin{emeraldbox}[title=Use Case Diagram (UCD)]
A use case diagram is a type of UML diagram that represents functional requirements. It visually describes interactions between users (actors) and the system, clarifying the operations and actions they can perform.
\end{emeraldbox}

\subsection{Actor Identification}

\begin{table}[H]
  \centering
  \caption{System actors and their roles}
  \label{tab:cs_actors}
  \renewcommand{\arraystretch}{1.3}
  \begin{tabular}{lp{10cm}}
    \toprule
    \textbf{Actor} & \textbf{Role Description} \\
    \midrule
    Visitor & An unauthenticated user who can sign up or recover a forgotten password. \\
    Learner & A registered user who practices Quran recitation, receives AI and teacher feedback, and tracks progress. \\
    Teacher & A qualified instructor who reviews student recitations, provides manual feedback, manages students, and issues Ijaza certificates. \\
    Admin & A system supervisor with full control over the platform, managing users, teachers, analytics, daily verse, and AI health. \\
    Tahkik AI Server & The backend AI system (NeMo FastConformer model) that processes audio and generates automated feedback. \\
    Supabase & The backend-as-a-service handling authentication, database, and real-time subscriptions. \\
    \bottomrule
  \end{tabular}
\end{table}

\subsection{Global Use Case Diagram}
The following use case diagram illustrates the global functionalities of the Tahkik platform across all actors.

\begin{figure}[H]
  \centering
  \includegraphics[width=\textwidth,height=0.9\textheight,keepaspectratio]{figures/use_case_diagram.png}
  \caption{Global Use Case Diagram for the Tahkik Platform}
  \label{fig:cs_global_ucd}
\end{figure}
\begin{table}[H]
  \centering
  \caption{UCD of Authentication description}
  \label{tab:cs_ucd_auth}
  \renewcommand{\arraystretch}{1.3}
  \begin{tabular}{lp{3cm}p{7cm}}
    \toprule
    \textbf{Action} & \textbf{Actors} & \textbf{Description} \\
    \midrule
    Sign up & Visitor & User creates an account by filling name, email, password, and selecting preferences. \\
    Log in & Learner, Teacher & User logs in with email/password or via Google Sign-In. \\
    Forgot password & Visitor & Initiates a password recovery process via email. \\
    Log out & All users & User logs out and session is terminated. \\
    \bottomrule
  \end{tabular}
\end{table}

\subsection{Use Case Diagram: Quran Recitation and AI Feedback}
The following use case diagram illustrates the core recitation and AI feedback features.

\begin{figure}[H]
  \centering
  \includegraphics[width=\textwidth,height=0.8\textheight,keepaspectratio]{figures/ucd_recitation.png}
  \caption{Use case diagram for Quran Recitation and AI Feedback}
  \label{fig:cs_ucd_recitation}
\end{figure}

\begin{table}[H]
  \centering
  \caption{UCD of Quran Recitation description}
  \label{tab:cs_ucd_recitation}
  \renewcommand{\arraystretch}{1.3}
  \begin{tabular}{lp{3cm}p{7cm}}
    \toprule
    \textbf{Action} & \textbf{Actors} & \textbf{Description} \\
    \midrule
    Browse Quran & Learner & Learner navigates the Quran reader with Warsh font rendering. \\
    Record Recitation & Learner & Learner records audio of their voice reciting selected verses. \\
    Submit Audio & Learner & Audio is streamed via WebSocket to the AI Server for processing. \\
    Get AI Feedback & Tahkik AI & Whisper model transcribes audio, detects errors, returns confidence score. \\
    View Score & Learner & Learner views the AI-generated feedback with highlighted errors. \\
    \bottomrule
  \end{tabular}
\end{table}

\subsection{Use Case Diagram: Teacher Review and Student Management}

\begin{figure}[H]
  \centering
  \includegraphics[width=\textwidth,height=0.8\textheight,keepaspectratio]{figures/ucd_teacher.png}
  \caption{Use case diagram for Teacher Review and Student Management}
  \label{fig:cs_ucd_teacher}
\end{figure}

\begin{table}[H]
  \centering
  \caption{UCD of Teacher Review description}
  \label{tab:cs_ucd_teacher}
  \renewcommand{\arraystretch}{1.3}
  \begin{tabular}{lp{3cm}p{7cm}}
    \toprule
    \textbf{Action} & \textbf{Actors} & \textbf{Description} \\
    \midrule
    Review Recitations & Teacher & Teacher views a queue of pending student recitations with AI flags. \\
    Add Manual Feedback & Teacher & Teacher listens to audio, reviews AI output, and adds corrections. \\
    Manage Students & Teacher & Teacher views student list, progress, and individual details. \\
    Issue Ijaza & Teacher & Teacher issues a digital Ijaza certificate to qualifying students. \\
    Send Messages & Teacher & Teacher communicates with students through in-app messaging. \\
    \bottomrule
  \end{tabular}
\end{table}

\subsection{Use Case Diagram: Administration}

\begin{figure}[H]
  \centering
  \includegraphics[width=\textwidth,height=0.8\textheight,keepaspectratio]{figures/ucd_admin.png}
  \caption{Use case diagram for Administration}
  \label{fig:cs_ucd_admin}
\end{figure}

\begin{table}[H]
  \centering
  \caption{UCD of Administration description}
  \label{tab:cs_ucd_admin}
  \renewcommand{\arraystretch}{1.3}
  \begin{tabular}{lp{3cm}p{7cm}}
    \toprule
    \textbf{Action} & \textbf{Actors} & \textbf{Description} \\
    \midrule
    Manage Users & Admin & View, search and update all registered users and their roles. \\
    Assign Roles & Admin & Promote or demote users between Learner, Teacher, and Admin roles. \\
    Manage Teachers & Admin & Promote users to teachers, assign students, and manage teacher rosters. \\
    Manage Subscriptions & Admin & View and manage subscription tiers and user distribution. \\
    View Platform Analytics & Admin & Access engagement metrics, usage trends, and recitation statistics. \\
    Monitor AI Health & Admin & View live inference server telemetry (CPU, memory, latency, container status). \\
    Manage Daily Verse & Admin & Curate and publish the daily verse displayed on the mobile home screen. \\
    \bottomrule
  \end{tabular}
\end{table}

\subsection{Use Case Diagram: Progress Tracking and Subscription}

\begin{figure}[H]
  \centering
  \includegraphics[width=\textwidth,height=0.8\textheight,keepaspectratio]{figures/ucd_progress.png}
  \caption{Use case diagram for Progress Tracking and Subscription}
  \label{fig:cs_ucd_progress}
\end{figure}

\begin{table}[H]
  \centering
  \caption{UCD of Progress Tracking and Subscription description}
  \label{tab:cs_ucd_progress}
  \renewcommand{\arraystretch}{1.3}
  \begin{tabular}{lp{3cm}p{7cm}}
    \toprule
    \textbf{Action} & \textbf{Actors} & \textbf{Description} \\
    \midrule
    View Progress Dashboard & Learner & Access charts, streak calendar, and aggregated performance data. \\
    View Recitation History & Learner & Browse past recitation sessions with scores and audio replay. \\
    View Streak & Learner & View current and longest streak of consecutive practice days. \\
    Track Daily Goal & Learner & Monitor daily practice minutes against the user-defined goal. \\
    View Average Scores & Learner & View average confidence scores per surah over time. \\
    View Ijaza Certificates & Learner & View digital Ijaza certificates issued by teachers. \\
    Manage Subscription & Learner & View current plan and available tiers. \\
    Upgrade to Premium & Learner & Subscribe to the premium tier for advanced features. \\
    \bottomrule
  \end{tabular}
\end{table}

\section{Sequence Diagrams}
\label{sec:cs_sequence}

\begin{emeraldbox}[title=Sequence Diagram]
UML sequence diagrams model the interactions and behaviors of system components over time. They visually represent how objects communicate through messages, making them valuable for requirements clarification and system design.
\end{emeraldbox}

\subsection{New Actors and Services Identification}

\begin{table}[H]
  \centering
  \caption{Actors of sequence diagrams}
  \label{tab:cs_seq_actors}
  \renewcommand{\arraystretch}{1.3}
  \begin{tabular}{lp{10cm}}
    \toprule
    \textbf{Actor} & \textbf{Role} \\
    \midrule
    Tahkik App (Flutter) & The mobile client application used by learners. \\
    Tahkik Dashboard (React) & The web dashboard used by teachers and admins. \\
    Supabase Auth & Authentication service handling JWT tokens and session management. \\
    Supabase DB & PostgreSQL database storing all application data. \\
    Tahkik AI Server & Backend server (Python/Whisper) that processes audio and returns AI feedback. \\
    \bottomrule
  \end{tabular}
\end{table}

\subsection{Sequence Diagram: Login}
The following sequence diagram illustrates the login process for a Learner using the Tahkik mobile app.

\begin{figure}[H]
  \centering
  \includegraphics[width=\textwidth,height=0.8\textheight,keepaspectratio]{figures/sequence_login.png}
  \caption{Sequence diagram for Login}
  \label{fig:cs_seq_login}
\end{figure}

\subsection{Sequence Diagram: Sign-up}
The following sequence diagram illustrates the sign-up process.

\begin{figure}[H]
  \centering
  \includegraphics[width=\textwidth,height=0.8\textheight,keepaspectratio]{figures/seq_signup.png}
  \caption{Sequence diagram for Sign-up}
  \label{fig:cs_seq_signup}
\end{figure}

\subsection{Sequence Diagram: Recitation Submission and AI Feedback}
The following sequence diagram illustrates the core recitation and AI feedback process.

\begin{figure}[H]
  \centering
  \includegraphics[width=\textwidth,height=0.8\textheight,keepaspectratio]{figures/sequence_recitation.png}
  \caption{Sequence diagram for Recitation and AI Feedback}
  \label{fig:cs_seq_recitation}
\end{figure}

\subsection{Sequence Diagram: Teacher Recitation Review}
The following sequence diagram illustrates the teacher review process.

\begin{figure}[H]
  \centering
  \includegraphics[width=\textwidth,height=0.8\textheight,keepaspectratio]{figures/seq_teacher.png}
  \caption{Sequence diagram for Teacher Recitation Review}
  \label{fig:cs_seq_teacher_review}
\end{figure}

\subsection{Sequence Diagram: Progress Tracking}
The following sequence diagram illustrates the progress tracking process.

\begin{figure}[H]
  \centering
  \includegraphics[width=\textwidth,height=0.8\textheight,keepaspectratio]{figures/seq_progress.png}
  \caption{Sequence diagram for Progress Tracking}
  \label{fig:cs_seq_progress}
\end{figure}

\subsection{Sequence Diagram: Admin User Management}
The following sequence diagram illustrates the admin management process.

\begin{figure}[H]
  \centering
  \includegraphics[width=\textwidth,height=0.8\textheight,keepaspectratio]{figures/seq_admin.png}
  \caption{Sequence diagram for Admin User Management}
  \label{fig:cs_seq_admin}
\end{figure}

\section{Class Diagram}
\label{sec:cs_class}

\begin{emeraldbox}[title=Class Diagram]
A class diagram is a crucial UML diagram representing both static and behavioral structures of an object-oriented system. It shows the classes, their attributes, operations, and the relationships among them.
\end{emeraldbox}

\subsection{Class Diagram Presentation}

\begin{figure}[H]
  \centering
  \includegraphics[width=\textwidth,height=0.85\textheight,keepaspectratio]{figures/class_diagram.png}
  \caption{Class diagram of the Tahkik platform}
  \label{fig:cs_class_diagram}
\end{figure}

\subsection{Class Diagram Description}

\begin{table}[H]
  \centering
  \caption{Description of Class diagram entities}
  \label{tab:cs_class_desc}
  \renewcommand{\arraystretch}{1.3}
  \begin{tabular}{lp{10cm}}
    \toprule
    \textbf{Class} & \textbf{Description} \\
    \midrule
    User & Abstract base class for all system users with common attributes (id, email, name, role). \\
    Learner & Inherits from User. Represents learners who practice recitation and track progress. \\
    Teacher & Inherits from User. Represents teachers who review recitations and manage students. \\
    Admin & Inherits from User. Represents administrators who manage the platform. \\
    Recitation & A recorded recitation submission with surah, ayah range, and audio URL. \\
    Feedback & Feedback on a recitation, either AI-generated or manual (from teacher), with confidence score and notes. \\
    Subscription & Tracks user subscription plans (free, premium) with start and end dates. \\
    Progress & Aggregates a learner's performance data (total recitations, average score, streak). \\
    Ijaza & A digital certification issued by a teacher to a qualifying learner for a specific surah. \\
    Message & Messages exchanged between users, with sender and receiver foreign keys. \\
    DailyVerse & A curated Quranic verse managed by administrators and displayed on the mobile home screen. \\
    TajweedResult & Per-word Tajweed analysis result for a recitation, storing rule type, severity, and timing. \\
    \bottomrule
  \end{tabular}
\end{table}

\section{Database Design}
\label{sec:cs_database}

\subsection{Relational Schema}
The relational model was derived from the UML class diagram by mapping each class into a relation:

\begin{infobox}[title=Relational Schema]
\textbf{User}(\underline{id}, email, name, role, avatar\_url, created\_at)

\textbf{Recitation}(\underline{id}, \textit{user\_id}, surah, ayah\_start, ayah\_end, audio\_url, submitted\_at)

\textbf{Feedback}(\underline{id}, \textit{recitation\_id}, \textit{teacher\_id}, type, confidence\_score, notes, created\_at)

\textbf{Subscription}(\underline{id}, \textit{user\_id}, plan\_type, start\_date, end\_date)

\textbf{Progress}(\underline{id}, \textit{user\_id}, total\_recitations, average\_score, streak\_days)

\textbf{Ijaza}(\underline{id}, \textit{learner\_id}, \textit{teacher\_id}, surah, issued\_at)

\textbf{Message}(\underline{id}, \textit{sender\_id}, \textit{receiver\_id}, content, sent\_at)

\textbf{DailyVerse}(\underline{id}, arabic\_text, translation, surah\_name, ayah\_number, is\_active, created\_at)

\textbf{TajweedResult}(\underline{id}, \textit{recitation\_id}, word\_index, rule\_type, severity, start\_time, end\_time)
\end{infobox}

Where \underline{underlined} attributes are primary keys and \textit{italicized} attributes are foreign keys.

\section{System Architecture}
\label{sec:cs_architecture}

\begin{figure}[H]
  \centering
  \includegraphics[width=\textwidth,height=0.8\textheight,keepaspectratio]{figures/sys_arch.png}
  \caption{System architecture diagram}
  \label{fig:cs_architecture}
\end{figure}

\section{Activity Diagram}
\label{sec:cs_activity}

\begin{emeraldbox}[title=Activity Diagram]
An activity diagram is a UML behavioral diagram that models the dynamic workflow of a system. It illustrates the sequence of activities, decisions, and parallel processes involved in a specific use case, providing a clear view of the control flow from start to finish.
\end{emeraldbox}

\subsection{Activity Diagram: Recitation and AI Feedback Flow}
The following activity diagram illustrates the complete workflow of the core recitation feature, from the moment the learner opens the Quran reader to the final progress update.

\begin{figure}[H]
  \centering
  \includegraphics[width=0.7\textwidth,height=0.85\textheight,keepaspectratio]{figures/activity_recitation.png}
  \caption{Activity diagram for the Recitation and AI Feedback flow}
  \label{fig:cs_activity_recitation}
\end{figure}

\section{Deployment Diagram}
\label{sec:cs_deployment}

\begin{emeraldbox}[title=Deployment Diagram]
A deployment diagram is a UML structural diagram that models the physical deployment of software artifacts onto hardware nodes. It shows the runtime architecture of the system, including servers, containers, and communication protocols.
\end{emeraldbox}

\subsection{Deployment Diagram Presentation}
The Tahkik platform is deployed across four environments: a mobile device (iOS/Android) running the Flutter app, a desktop browser hosting the React dashboard, a DigitalOcean droplet running a Docker Compose stack (Caddy reverse proxy, Go API gateway, and Python NeMo inference service), and the Supabase cloud providing authentication, database, storage, and real-time services. Model weights are fetched from the Hugging Face Hub.

\begin{figure}[H]
  \centering
  \includegraphics[width=\textwidth,height=0.85\textheight,keepaspectratio]{figures/deployment_diagram.png}
  \caption{Deployment diagram of the Tahkik platform}
  \label{fig:cs_deployment}
\end{figure}

\section{Conclusion}
This chapter presented the conceptual foundation of the Tahkik platform. We identified 10 functional and 5 non-functional requirements, designed the system's behavioral interactions through 5 use case diagrams, 6 sequence diagrams, and 1 activity diagram, and defined the static structure via a class diagram, relational schema, system architecture, and deployment diagram. The next chapter details the implementation and realization of this design.
